Noisy distances from a point to known anchors give posteriors whose credible regions are rings, shells, tubes and
disconnected pieces rather than ellipses. We determine their topology exactly. The key observation is that for two ranges,
in residual coordinates (the differences between distances and observed ranges), the posterior is a Gaussian times an
explicit weight on a convex polygonal domain, and every highest-posterior-density region is a disc cut by the walls
of that domain. Its topology can therefore be read off the triangle whose sides are the anchor separation and the two
observed ranges. This gives the complete sequence of topologies as the noise grows, in every dimension: each transition
occurs when the cut-off reaches one of five depths, the squared incircle tangent lengths $t_*^2, t_1^2, t_2^2$ and the
sums $t_*^2 + t_1^2$, $t_*^2 + t_2^2$, and more noise first creates holes and then removes them. With a different noise
level on each range the disc becomes an ellipse, and a contact can annihilate with an anchor. For $n \le d$ ranges the
noise-free shape persists while the region misses the affine hull of the anchors, and for $n = d$ only then. An explicit
noise level, computable from the data, certifies the noise-free shape. In simulation, from the noise level at which the first hole forms, even a mixture of
Laplace approximations at the two posterior modes covers the true point in only 84--93\% of data sets at the 95\% level. For one range the region has a hole exactly
when $r/\sigma$ exceeds $\kappa_d = \sqrt{d/2} + z\sqrt{3/8} + O(d^{-1/2})$, and the posterior mean then lies in the hole.
In high dimension every transition has an explicit limit with a $1/\sqrt d$ correction, and for random anchors in the
plane the flip pair is lost with probability of order $\sqrt\sigma$, up to a logarithmic factor.
